\documentclass[10pt]{article}
\usepackage[margin=0.8in]{geometry}
\usepackage[]{xcolor}
\usepackage{amsmath}
\parindent=0pt

\title{CSE 4300: HW 3}
\author{Gordon Chen, gbc24001}
\date{}

\begin{document}
\maketitle


\section*{Completed Parts}
I completed the all parts of the hw. I added the \verb|_exit|, \verb|printint|, and \verb|reversestring|
syscalls, added a test for all three syscalls, and answered the questions for Part E.



\section*{Changes}
\subsection*{syscall.h}
\begin{verbatim}
void sys__exit(int exitCode);

int sys_printint(int c);

int sys_reversestring(const char *str, int len);
\end{verbatim}

Add kernel-level syscall function prototypes.


\subsection*{simple\_syscalls.c}
\begin{verbatim}
#include <types.h>
#include <syscall.h>
#include <thread.h>
#include <lib.h>

void sys__exit(int exitCode) {
    kprintf("_exit exitCode: %d\n", exitCode);
    thread_exit(exitCode);
    panic("_exit failed to kill thread\n");
}

int sys_printint(int c) {
    kprintf("%d\n", c);
    return (c % 5 != 0);
}

int sys_reversestring(const char *str, int len) {
    int og_len = len;
    --len;
    while (len >= 0) {
        kprintf("%c", *(str+len));
        --len;
    }
    kprintf("\n");
    return (og_len % 2 == 0);
}
\end{verbatim}

Add syscall implementations.


\subsection*{thread.c}
\begin{verbatim}
void thread_exit(int exitCode) {
    kprintf("thread_exit exitCode: %d\n", exitCode);
\end{verbatim}

update \verb|thread_exit| to take \verb|exitCode|. Also changed all \verb|thread_exit| calls to take
the new \verb|exitCode| parameter.


\subsection*{conf.kern}
\begin{verbatim}
file      userprog/simple_syscalls.c
\end{verbatim}

Add \verb|simple_syscalls| to kern conf.


\subsection*{callno.h}
\begin{verbatim}
#define SYS_printint     32
#define SYS_reversestring 33
\end{verbatim}

Define syscall numbers for \verb|printint| and \verb|reversestring|.


\subsection*{syscall.c}
\begin{verbatim}
int err = 0;
...
case SYS__exit:
    sys__exit(tf->tf_a0);
    break;

case SYS_printint:
    retval = sys_printint(tf->tf_a0);
    break;

case SYS_reversestring:
    retval = sys_reversestring(tf->tf_a0, tf->tf_a1);
    break;
\end{verbatim}

Init \verb|err| to 0. Match syscall number and call syscall function.


\subsection*{unistd.h}
\begin{verbatim}
int printint(int c);
int reversestring(const char *str, int len);
\end{verbatim}

Add user-level function prototypes to \verb|unistd.h|.


\section*{testA}
\subsection*{testA.c}
\begin{verbatim}
int main() {
    int ret = printint(5);
    printint(ret);  // 0 (divisible by 5).

    ret = printint(6);
    printint(ret);  // 1.

    ret = reversestring("hello", 5);
    printint(ret);  // 0 (not divisible by 2).

    ret = reversestring("hell", 4);
    printint(ret);  // 1.

    _exit(0);
    printint(-1);  // This should not print.
}
\end{verbatim}

\subsection*{testA output}
\begin{verbatim}
OS/161 kernel [? for menu]: p testbin/testA
5
0
6
1
olleh
0
lleh
1
_exit exitCode: 0
thread_exit exitCode: 0
Operation took 0.160804800 seconds
\end{verbatim}

The output is correct. \verb|printint| first prints 5. Then returns 0 since 5 is divisible by 5.
Then \verb|printint| prints 6, and returns 1 since 6 is not divisible by 5.
Then \verb|reversestring| prints \verb|olleh| and returns 0 since 5 is not divisible by 2.
Then \verb|reversestring| prints \verb|lleh| and returns 1 since 4 is divisible by 2.
Then we \verb|_exit| and see that the exit code is printed in both \verb|_exit| and \verb|thread_exit|.
\verb|printint| never prints -1 because we exited.


\section*{Part E Answers}
\begin{enumerate}
    \item How do you enable and disable interrupt in OS161?

        enable: call \verb|spl0|, which calls \verb|splx|, which calls \verb|interrupts_on|,
            which sets the \verb|IEc| (global interrupt enable) flag.

        disable: call \verb|splhigh|, which calls \verb|splx|, which calls \verb|interrupts_off|,
            which clears the \verb|IEc| (global interrupt enable) flag.

    \item What is the purpose of thread\_yield() function?

        To yield the CPU to another thread but stay runnable.

    \item What is the difference between thread\_yield() and thread\_sleep() function?

        \verb|thread_yield| yields to another thread but stays runnable.
        \verb|thread_sleep| yield to another thread but is not immediately runnable,
        it must be woken up with \verb|thread_wakeup|.

    \item What is the purpose of mi\_switch function?

        To perform a context switch. Checks magic numbers at the bottom of the stack to ensure that the
        function didn't overflow the stack. Then adds the thread calling the switch to the runnable or
        sleeping thread queue depending on if \verb|mi_switch| was called by by \verb|thread_yield|
        or \verb|thread_sleep| respectively. Then calls the scheduler to get the next thread to run,
        and calls a function to perform the actual context switch.

    \item Where is the scheduler function located and what does the current scheduler do?

        The scheduler function is in \verb|kern/thread/scheduler.c|. The current scheduler
        waits for a thread to enter the runnable queue, and then returns the head of the runnable queue,
        ie the next thread to run. So the scheduler runs threads in a FIFO queue, it uses
        round-robin scheduling.

\end{enumerate}


\end{document}
